\documentclass[aspectratio=169]{beamer}
\usepackage[utf8]{inputenc}
\usepackage{amsmath}
\usepackage{amssymb}
\usepackage{listings}
\usepackage{tikz}
\usepackage{booktabs}
\usepackage{multicol}
\usepackage{hyperref}

\usetikzlibrary{shapes,arrows,positioning,fit,backgrounds}

% Theme settings
% \usetheme{Madrid}
% \usecolortheme{default}
% \setbeamertemplate{navigation symbols}{}
% \setbeamertemplate{footline}[frame number]
 %% Beamer Layout %%%%%%%%%%%%%%%%%%%%%%%%%%%%%%%%%%
 %\useoutertheme[subsection=false,shadow]{miniframes}
 %\useinnertheme{default}
 % \usecolortheme{seahorse}
 \usepackage{beamerthemesplit} 
 \setbeamercolor{math text}{fg=blue!80!black}
 \setbeamercolor{normal text}{fg=blue!80!blue}
 \usetheme{AnnArbor}
 %\useinnertheme[subsection=false,shadow]{rounded}
 %\useoutertheme{infolines}
 \usecolortheme{seahorse}
 \usefonttheme{serif,professionalfonts}
 %\usefonttheme{platino}

 
 \setbeamercolor{header}{fg=white,bg=DeepSkyBlue4}
 \setbeamercolor{frametitle}{bg=white,fg=DeepSkyBlue4}%, fontface=bf}
 \setbeamercolor*{section}{fg=white,bg=DeepSkyBlue4}%, fontface=bf}
 \setbeamercolor{title}{fg=white,bg=DeepSkyBlue4}
 \setbeamercolor{item}{fg=DeepSkyBlue4}
 \setbeamercolor*{lower separation line head}{bg=DeepSkyBlue4} 
 \setbeamercolor*{normal text}{fg=black,bg=white} 
 \setbeamercolor*{alerted text}{fg=red} 
 \setbeamercolor*{example text}{fg=black} 
 \setbeamercolor*{structure}{fg=black} 
 \setbeamercolor*{math text}{fg=blue!80!black}
 \setbeamercolor*{palette tertiary}{fg=black,bg=black!10} 
 \setbeamercolor*{palette quaternary}{fg=black,bg=black!10} 

% Color definitions
\definecolor{googleblue}{RGB}{66,133,244}
\definecolor{googlered}{RGB}{219,68,55}
\definecolor{googleyellow}{RGB}{244,160,0}
\definecolor{googlegreen}{RGB}{15,157,88}
\definecolor{codebg}{RGB}{245,245,245}

\setbeamercolor{structure}{fg=googleblue}
\setbeamercolor{title}{fg=googleblue}
\setbeamercolor{frametitle}{fg=googleblue}

% Code listing settings
\lstset{
    basicstyle=\ttfamily\footnotesize,
    backgroundcolor=\color{codebg},
    frame=single,
    breaklines=true,
    columns=flexible,
    keywordstyle=\color{googleblue}\bfseries,
    commentstyle=\color{googlegreen}\itshape,
    stringstyle=\color{googlered},
    numbers=left,
    numberstyle=\tiny\color{gray},
    showstringspaces=false,
}

% Title page
\title[MapReduce Tutorial]{MapReduce: Simplified Data Processing \\ on Large Clusters}
\subtitle{A Tutorial Based on Google's Seminal Paper}
\author{Jeffrey Dean and Sanjay Ghemawat}
\institute{Google, Inc. \\ \vspace{0.3cm} \textit{Presented at OSDI 2004}}
\date{}

\begin{document}

% ============================================
% TITLE SLIDE
% ============================================
\begin{frame}
\maketitle
\begin{center}
    \vspace{-0.5cm}
    \includegraphics[width=0.15\textwidth]{example-image} % Replace with Google logo if available
    \vspace{0.3cm}

    \footnotesize
    Tutorial based on the paper: \\
    \textit{"MapReduce: Simplified Data Processing on Large Clusters"} \\
    OSDI 2004
\end{center}
\end{frame}

% ============================================
% TABLE OF CONTENTS
% ============================================
\begin{frame}{Outline}
\tableofcontents
\end{frame}

% ============================================
% SECTION 1: INTRODUCTION
% ============================================
\section{Introduction and Motivation}

\begin{frame}{The Problem at Google (circa 2004)}
\begin{columns}
\column{0.5\textwidth}
\textbf{Challenges:}
\begin{itemize}
    \item Processing \alert{terabytes} of data
    \item Hundreds of special-purpose computations
    \item Examples:
    \begin{itemize}
        \item Crawled documents
        \item Web request logs
        \item Inverted indices
        \item Graph structures
        \item Page summaries
    \end{itemize}
\end{itemize}

\column{0.5\textwidth}
\textbf{Issues:}
\begin{itemize}
    \item Input data is \alert{huge}
    \item Must distribute across \alert{hundreds/thousands} of machines
    \item Complex code for:
    \begin{itemize}
        \item Parallelization
        \item Data distribution
        \item Fault tolerance
        \item Load balancing
    \end{itemize}
\end{itemize}
\end{columns}

\vspace{0.5cm}
\begin{alertblock}{Core Problem}
Simple computations obscured by large amounts of complex code dealing with distribution and fault tolerance
\end{alertblock}
\end{frame}

\begin{frame}{The MapReduce Solution}
\begin{block}{Key Insight}
Most computations involve:
\begin{enumerate}
    \item Applying a \textbf{map} operation to each logical record
    \item Applying a \textbf{reduce} operation to all values sharing the same key
\end{enumerate}
\end{block}

\vspace{0.3cm}

\textbf{Inspiration:} Map and reduce primitives from functional programming (Lisp, etc.)

\vspace{0.3cm}

\begin{columns}
\column{0.5\textwidth}
\textbf{What MapReduce Provides:}
\begin{itemize}
    \item Simple programming interface
    \item Automatic parallelization
    \item Fault tolerance via re-execution
    \item Load balancing
\end{itemize}

\column{0.5\textwidth}
\textbf{Benefits:}
\begin{itemize}
    \item Hides messy details
    \item Programmers focus on logic
    \item No distributed systems expertise needed
    \item Scales to thousands of machines
\end{itemize}
\end{columns}
\end{frame}

\begin{frame}{Major Contributions}
\begin{enumerate}
    \item \textbf{Simple and Powerful Interface}
    \begin{itemize}
        \item Enables automatic parallelization and distribution
        \item Easy to express large-scale computations
    \end{itemize}

    \vspace{0.3cm}

    \item \textbf{High-Performance Implementation}
    \begin{itemize}
        \item Achieves high performance on large clusters of commodity PCs
        \item Processes many terabytes on thousands of machines
    \end{itemize}

    \vspace{0.3cm}

    \item \textbf{Proven at Scale}
    \begin{itemize}
        \item Hundreds of MapReduce programs implemented at Google
        \item Thousands of jobs executed daily on Google's clusters
    \end{itemize}
\end{enumerate}
\end{frame}

% ============================================
% SECTION 2: PROGRAMMING MODEL
% ============================================
\section{Programming Model}

\begin{frame}{MapReduce Programming Model}
\begin{block}{Core Concept}
Computation takes a set of \textbf{input} key/value pairs and produces a set of \textbf{output} key/value pairs
\end{block}

\vspace{0.3cm}

\textbf{User specifies two functions:}

\begin{enumerate}
    \item \textbf{Map Function}
    \begin{itemize}
        \item Takes an input pair
        \item Produces a set of \textit{intermediate} key/value pairs
        \item Library groups all intermediate values by the same intermediate key
    \end{itemize}

    \vspace{0.2cm}

    \item \textbf{Reduce Function}
    \begin{itemize}
        \item Accepts an intermediate key and a set of values for that key
        \item Merges values to form a possibly smaller set
        \item Typically produces zero or one output value per invocation
    \end{itemize}
\end{enumerate}
\end{frame}

\begin{frame}{Type Signatures}
\begin{block}{Conceptual Types}
Even though implementations may use strings, conceptually:

\vspace{0.3cm}

\centering
\begin{tabular}{ll}
\textbf{Map:} & $(k_1, v_1) \rightarrow \text{list}(k_2, v_2)$ \\[0.3cm]
\textbf{Reduce:} & $(k_2, \text{list}(v_2)) \rightarrow \text{list}(v_2)$ \\
\end{tabular}
\end{block}

\vspace{0.5cm}

\textbf{Important Notes:}
\begin{itemize}
    \item Input keys/values ($k_1, v_1$) are from a \alert{different domain} than output keys/values
    \item Intermediate keys/values ($k_2, v_2$) are from the \alert{same domain} as output keys/values
    \item Values supplied to Reduce via an \textbf{iterator} (handles large lists)
\end{itemize}
\end{frame}

\begin{frame}[fragile]{Classic Example: Word Count}
\textbf{Problem:} Count occurrences of each word in a large collection of documents

\vspace{0.3cm}

\begin{lstlisting}[language=Python, caption=Map Function]
map(String key, String value):
    // key: document name
    // value: document contents
    for each word w in value:
        EmitIntermediate(w, "1");
\end{lstlisting}

\vspace{0.2cm}

\begin{lstlisting}[language=Python, caption=Reduce Function]
reduce(String key, Iterator values):
    // key: a word
    // values: a list of counts
    int result = 0;
    for each v in values:
        result += ParseInt(v);
    Emit(AsString(result));
\end{lstlisting}
\end{frame}

\begin{frame}{Word Count: Execution Flow}
\begin{center}
\begin{tikzpicture}[
    node distance=1.5cm,
    box/.style={rectangle, draw, minimum width=2cm, minimum height=0.8cm, align=center},
    arrow/.style={->, >=stealth, thick}
]

% Input
\node[box, fill=googleblue!20] (input) {Input Documents};

% Map phase
\node[box, fill=googleyellow!20, below of=input] (map1) {Map 1};
\node[box, fill=googleyellow!20, right=0.5cm of map1] (map2) {Map 2};
\node[box, fill=googleyellow!20, right=0.5cm of map2] (map3) {Map 3};

% Intermediate
\node[box, fill=googlegreen!20, below of=map2] (shuffle) {Shuffle \& Sort \\ by Key};

% Reduce phase
\node[box, fill=googlered!20, below of=shuffle, xshift=-1.5cm] (reduce1) {Reduce 1};
\node[box, fill=googlered!20, below of=shuffle, xshift=1.5cm] (reduce2) {Reduce 2};

% Output
\node[box, fill=googleblue!20, below of=shuffle, yshift=-2cm] (output) {Output Files};

% Arrows
\draw[arrow] (input) -- (map1);
\draw[arrow] (input) -- (map2);
\draw[arrow] (input) -- (map3);

\draw[arrow] (map1) -- (shuffle);
\draw[arrow] (map2) -- (shuffle);
\draw[arrow] (map3) -- (shuffle);

\draw[arrow] (shuffle) -- (reduce1);
\draw[arrow] (shuffle) -- (reduce2);

\draw[arrow] (reduce1) -- (output);
\draw[arrow] (reduce2) -- (output);

% Annotations
\node[right=0.2cm of map3, text width=2.5cm, font=\tiny] {Emit: (word, "1") for each word};
\node[right=0.2cm of reduce2, text width=2.5cm, font=\tiny] {Emit: (word, count)};

\end{tikzpicture}
\end{center}
\end{frame}

\begin{frame}{More Examples}
\begin{block}{Distributed Grep}
\textbf{Map:} Emit line if it matches pattern \\
\textbf{Reduce:} Identity function (copies data to output)
\end{block}

\begin{block}{Count of URL Access Frequency}
\textbf{Map:} Process web logs, output $\langle\text{URL}, 1\rangle$ \\
\textbf{Reduce:} Add all values for same URL, emit $\langle\text{URL}, \text{total count}\rangle$
\end{block}

\begin{block}{Reverse Web-Link Graph}
\textbf{Map:} For each link to target in page source, output $\langle\text{target}, \text{source}\rangle$ \\
\textbf{Reduce:} Concatenate all sources, emit $\langle\text{target}, \text{list(source)}\rangle$
\end{block}

\begin{block}{Inverted Index}
\textbf{Map:} Parse document, emit $\langle\text{word}, \text{document ID}\rangle$ \\
\textbf{Reduce:} Sort document IDs, emit $\langle\text{word}, \text{list(document ID)}\rangle$
\end{block}
\end{frame}

% ============================================
% SECTION 3: IMPLEMENTATION
% ============================================
\section{Implementation}

\begin{frame}{Google's Computing Environment}
\textbf{Target Platform:} Large clusters of commodity PCs connected by switched Ethernet

\vspace{0.3cm}

\begin{columns}
\column{0.5\textwidth}
\textbf{Hardware Specs:}
\begin{itemize}
    \item Dual-processor x86 (2GHz Xeon)
    \item HyperThreading enabled
    \item 2-4 GB memory per machine
    \item Two 160GB IDE disks
    \item Gigabit Ethernet
\end{itemize}

\column{0.5\textwidth}
\textbf{Environment:}
\begin{itemize}
    \item Hundreds/thousands of machines
    \item Two-level tree-switched network
    \item 100-200 Gbps aggregate bandwidth
    \item Machine failures are \alert{common}
    \item Distributed file system (GFS)
    \item Job scheduling system
\end{itemize}
\end{columns}

\vspace{0.5cm}

\begin{alertblock}{Key Insight}
Network bandwidth is a \textbf{scarce resource} $\rightarrow$ Optimize for locality!
\end{alertblock}
\end{frame}

\begin{frame}[fragile]{Execution Overview}
\begin{center}
\begin{tikzpicture}[
    node distance=0.8cm,
    box/.style={rectangle, draw, minimum width=1.2cm, minimum height=0.6cm, font=\tiny},
    process/.style={ellipse, draw, minimum width=1cm, minimum height=0.6cm, font=\tiny},
    arrow/.style={->, >=stealth}
]

% User program and master
\node[process, fill=googleblue!20] (user) {User Program};
\node[process, fill=googlered!20, below=1.5cm of user] (master) {Master};

% Workers
\node[process, fill=googleyellow!20, below left=1.5cm and 2cm of master] (worker1) {worker};
\node[process, fill=googleyellow!20, below=2.5cm of master] (worker2) {worker};
\node[process, fill=googleyellow!20, below right=1.5cm and 2cm of master] (worker3) {worker};
\node[process, fill=googlegreen!20, right=3cm of worker2] (worker4) {worker};
\node[process, fill=googlegreen!20, above=0.8cm of worker4] (worker5) {worker};

% Input splits
\node[box, fill=gray!20, left=1.5cm of worker1] (split0) {split 0};
\node[box, fill=gray!20, below=0.2cm of split0] (split1) {split 1};
\node[box, fill=gray!20, below=0.2cm of split1] (split2) {split 2};
\node[box, fill=gray!20, below=0.2cm of split2] (split3) {split 3};
\node[box, fill=gray!20, below=0.2cm of split3] (split4) {split 4};

% Intermediate files
\node[box, fill=orange!20, right=0.5cm of worker1, yshift=-0.5cm] (int1) {};
\node[box, fill=orange!20, right=0.5cm of worker2, yshift=-0.5cm] (int2) {};
\node[box, fill=orange!20, right=0.5cm of worker3, yshift=-0.5cm] (int3) {};

% Output files
\node[box, fill=googleblue!30, right=1.5cm of worker4] (out0) {output file 0};
\node[box, fill=googleblue!30, right=1.5cm of worker5] (out1) {output file 1};

% Labels
\node[above=0.1cm of split0, font=\tiny] {Input files};
\node[above right=0.1cm and -0.5cm of int2, font=\tiny] {Intermediate};
\node[above=0.1cm of out1, font=\tiny] {Output};

% Arrows with labels
\draw[arrow, dashed] (user) -- node[left, font=\tiny] {(1) fork} (master);
\draw[arrow, dashed] (user) -- (worker1);
\draw[arrow, dashed] (user) -- (worker3);

\draw[arrow] (master) -- node[above, font=\tiny] {(2) assign} (worker1);
\draw[arrow] (master) -- node[right, font=\tiny] {(2) assign} (worker4);

\draw[arrow] (split2) -- node[above, font=\tiny] {(3) read} (worker2);

\draw[arrow] (worker2) -- node[above, font=\tiny] {(4) local write} (int2);

\draw[arrow] (int2) -- node[above, font=\tiny] {(5) remote read} (worker4);

\draw[arrow] (worker4) -- node[above, font=\tiny] {(6) write} (out0);

% Phase labels
\node[below=2cm of split2, font=\tiny\bfseries, text=googleblue] {Map phase};
\node[below=2cm of worker4, font=\tiny\bfseries, text=googlegreen] {Reduce phase};

\end{tikzpicture}
\end{center}
\end{frame}

\begin{frame}{Execution Steps (1/2)}
\textbf{Step 1: Split Input}
\begin{itemize}
    \item MapReduce library splits input files into $M$ pieces (typically 16-64 MB)
    \item Starts up many copies of the program on cluster machines
\end{itemize}

\vspace{0.2cm}

\textbf{Step 2: Master Assignment}
\begin{itemize}
    \item One copy is special: the \alert{master}
    \item Rest are workers assigned work by the master
    \item Master assigns $M$ map tasks and $R$ reduce tasks to idle workers
\end{itemize}

\vspace{0.2cm}

\textbf{Step 3: Map Task Execution}
\begin{itemize}
    \item Worker reads contents of corresponding input split
    \item Parses key/value pairs from input data
    \item Passes each pair to user-defined Map function
    \item Intermediate key/value pairs buffered in memory
\end{itemize}
\end{frame}

\begin{frame}{Execution Steps (2/2)}
\textbf{Step 4: Write Intermediate Data}
\begin{itemize}
    \item Buffered pairs periodically written to local disk
    \item Partitioned into $R$ regions by partitioning function
    \item Locations passed back to master
\end{itemize}

\vspace{0.2cm}

\textbf{Step 5: Reduce Reads Intermediate Data}
\begin{itemize}
    \item Master notifies reduce worker about locations
    \item Reduce worker uses RPC to read buffered data from map workers' local disks
    \item Sorts intermediate data by key (groups occurrences of same key)
    \item External sort used if data too large for memory
\end{itemize}

\vspace{0.2cm}

\textbf{Steps 6-7: Reduce and Complete}
\begin{itemize}
    \item Reduce worker iterates over sorted data
    \item For each unique key, passes key and values to Reduce function
    \item Output appended to final output file for this reduce partition
    \item When all tasks complete, master wakes up user program
\end{itemize}
\end{frame}

\begin{frame}{Master Data Structures}
\textbf{State Tracking:}
\begin{itemize}
    \item For each map and reduce task:
    \begin{itemize}
        \item State: \textit{idle}, \textit{in-progress}, or \textit{completed}
        \item Identity of worker machine (for non-idle tasks)
    \end{itemize}
\end{itemize}

\vspace{0.3cm}

\textbf{Location Propagation:}
\begin{itemize}
    \item Master is conduit for intermediate file regions
    \item For each completed map task, stores:
    \begin{itemize}
        \item Locations of $R$ intermediate file regions
        \item Sizes of intermediate file regions
    \end{itemize}
    \item Information pushed incrementally to workers with in-progress reduce tasks
\end{itemize}

\vspace{0.3cm}

\begin{alertblock}{Memory Overhead}
$O(M + R)$ scheduling decisions and $O(M \times R)$ state in memory \\
(Constant factors are small: $\approx$ 1 byte per map/reduce task pair)
\end{alertblock}
\end{frame}

\begin{frame}{Fault Tolerance: Worker Failure}
\textbf{Detection:}
\begin{itemize}
    \item Master pings workers periodically
    \item No response in certain time $\rightarrow$ worker marked as failed
\end{itemize}

\vspace{0.3cm}

\textbf{Recovery:}
\begin{itemize}
    \item Completed map tasks reset to \textit{idle} state
    \item Map/reduce tasks in progress reset to \textit{idle}
    \item Eligible for rescheduling on other workers
\end{itemize}

\vspace{0.3cm}

\textbf{Why re-execute completed map tasks?}
\begin{itemize}
    \item Output stored on local disk of failed machine $\rightarrow$ \alert{inaccessible}
    \item Completed reduce tasks don't need re-execution
    \item Their output stored in global file system
\end{itemize}

\vspace{0.3cm}

\begin{exampleblock}{Resilience Example}
Network maintenance caused 80 machines to become unreachable. \\
MapReduce simply re-executed work and continued to completion!
\end{exampleblock}
\end{frame}

\begin{frame}{Fault Tolerance: Master Failure \& Semantics}
\textbf{Master Failure:}
\begin{itemize}
    \item Master writes periodic checkpoints
    \item New copy can be started from last checkpoint
    \item Current implementation: abort if master fails (clients can retry)
    \item Single master makes failure unlikely
\end{itemize}

\vspace{0.3cm}

\textbf{Semantics in Presence of Failures:}

\begin{block}{Deterministic Functions}
When map and reduce operators are \textbf{deterministic}:
\begin{itemize}
    \item Distributed implementation produces same output as non-faulting sequential execution
    \item Relies on atomic commits of map and reduce task outputs
\end{itemize}
\end{block}

\begin{block}{Non-Deterministic Functions}
Weaker but reasonable semantics:
\begin{itemize}
    \item Output equivalent to some sequential execution
    \item Different reduce tasks may see outputs from different executions of same map task
\end{itemize}
\end{block}
\end{frame}

\begin{frame}{Locality Optimization}
\begin{block}{Key Observation}
Network bandwidth is a \alert{relatively scarce resource}
\end{block}

\vspace{0.3cm}

\textbf{Strategy:}
\begin{itemize}
    \item Input data stored on local disks via GFS (Google File System)
    \item GFS divides files into 64 MB blocks
    \item Stores 3 copies of each block on different machines
    \item Master takes location information into account
\end{itemize}

\vspace{0.3cm}

\textbf{Scheduling Policy:}
\begin{enumerate}
    \item \textbf{Best:} Schedule map task on machine containing replica of input data
    \item \textbf{Next best:} Schedule near a replica (e.g., same network switch)
\end{enumerate}

\vspace{0.3cm}

\begin{exampleblock}{Result}
When running large MapReduce operations on significant fraction of workers: \\
\textbf{Most input data read locally} $\rightarrow$ consumes \textbf{no network bandwidth}
\end{exampleblock}
\end{frame}

\begin{frame}{Task Granularity}
\textbf{Subdividing:}
\begin{itemize}
    \item Map phase: $M$ pieces
    \item Reduce phase: $R$ pieces
\end{itemize}

\vspace{0.3cm}

\textbf{Ideally:} $M$ and $R$ should be much larger than number of worker machines

\vspace{0.3cm}

\textbf{Benefits:}
\begin{itemize}
    \item Improves dynamic load balancing
    \item Speeds up recovery when worker fails
    \item Completed tasks spread across all other workers
\end{itemize}

\vspace{0.3cm}

\textbf{Practical Bounds:}
\begin{itemize}
    \item Master makes $O(M + R)$ scheduling decisions
    \item Keeps $O(M \times R)$ state in memory
\end{itemize}

\vspace{0.3cm}

\begin{exampleblock}{Typical Configuration at Google}
$M = 200,000$, $R = 5,000$ using 2,000 worker machines \\
Each task: 16-64 MB of input data
\end{exampleblock}
\end{frame}

\begin{frame}{Backup Tasks: Fighting Stragglers}
\begin{block}{Problem: Stragglers}
A machine taking unusually long time to complete one of last few tasks
\end{block}

\vspace{0.3cm}

\textbf{Causes of Stragglers:}
\begin{itemize}
    \item Bad disk (30 MB/s $\rightarrow$ 1 MB/s due to correctable errors)
    \item Cluster scheduler scheduled other tasks (competition for resources)
    \item Bug in machine initialization (e.g., processor caches disabled)
\end{itemize}

\vspace{0.3cm}

\textbf{Solution: Backup Tasks}
\begin{itemize}
    \item When MapReduce operation close to completion
    \item Master schedules \textbf{backup executions} of remaining in-progress tasks
    \item Task marked complete when \alert{either} primary \alert{or} backup completes
    \item Tuned to increase resources by only a few percent
\end{itemize}

\vspace{0.3cm}

\begin{alertblock}{Impact}
Sort program: \textbf{44\% longer} to complete when backup tasks disabled!
\end{alertblock}
\end{frame}

% ============================================
% SECTION 4: REFINEMENTS
% ============================================
\section{Refinements}

\begin{frame}{Refinement 1: Partitioning Function}
\textbf{Default:} Hash-based partitioning
\begin{itemize}
    \item \texttt{hash(key) mod R}
    \item Results in fairly well-balanced partitions
\end{itemize}

\vspace{0.3cm}

\textbf{Custom Partitioning:}
\begin{itemize}
    \item Sometimes useful to partition by other functions
    \item \textbf{Example:} Output keys are URLs
    \item Want all URLs from same host in same output file
    \item Solution: \texttt{hash(Hostname(urlkey)) mod R}
\end{itemize}

\vspace{0.3cm}

\textbf{Benefits:}
\begin{itemize}
    \item Better data organization
    \item Enables efficient downstream processing
    \item Supports application-specific requirements
\end{itemize}
\end{frame}

\begin{frame}{Refinement 2: Ordering Guarantees}
\begin{block}{Guarantee}
Within a given partition, intermediate key/value pairs processed in \textbf{increasing key order}
\end{block}

\vspace{0.3cm}

\textbf{Benefits:}
\begin{itemize}
    \item Easy to generate sorted output file per partition
    \item Useful when output format needs efficient random access by key
    \item Convenient for users who want sorted data
\end{itemize}

\vspace{0.3cm}

\textbf{Use Cases:}
\begin{itemize}
    \item Building searchable indices
    \item Database-like access patterns
    \item Sorted data analysis
\end{itemize}
\end{frame}

\begin{frame}{Refinement 3: Combiner Function}
\begin{block}{Problem}
Significant repetition in intermediate keys produced by map tasks \\
\textbf{Example:} Word counting - many $\langle$the, 1$\rangle$ records
\end{block}

\vspace{0.3cm}

\textbf{Solution: Combiner Function}
\begin{itemize}
    \item Partial merging of data \alert{before} sending over network
    \item Executed on each machine performing map task
    \item Typically same code as reduce function
\end{itemize}

\vspace{0.3cm}

\textbf{Differences from Reduce:}
\begin{itemize}
    \item \textbf{Reduce:} Output written to final output file
    \item \textbf{Combiner:} Output written to intermediate file sent to reduce task
\end{itemize}

\vspace{0.3cm}

\begin{exampleblock}{Impact}
Partial combining \textbf{significantly speeds up} certain classes of MapReduce operations
\end{exampleblock}
\end{frame}

\begin{frame}{Other Refinements}
\begin{block}{4. Input and Output Types}
\begin{itemize}
    \item Support for reading data in several formats (text, binary, database, etc.)
    \item Users can add custom input/output types via reader interface
\end{itemize}
\end{block}

\begin{block}{5. Skipping Bad Records}
\begin{itemize}
    \item Detect records causing deterministic crashes
    \item Skip these records to make forward progress
    \item Useful when bugs in third-party libraries can't be fixed
\end{itemize}
\end{block}

\begin{block}{6. Local Execution}
\begin{itemize}
    \item Alternative implementation for debugging
    \item Executes all work sequentially on local machine
    \item Facilitates debugging, profiling, small-scale testing
\end{itemize}
\end{block}

\begin{block}{7. Status Information \& Counters}
\begin{itemize}
    \item HTTP server exports status pages (progress, failed workers, etc.)
    \item Counter facility for tracking events (words processed, documents indexed, etc.)
\end{itemize}
\end{block}
\end{frame}

% ============================================
% SECTION 5: PERFORMANCE
% ============================================
\section{Performance}

\begin{frame}{Cluster Configuration}
\textbf{Test Environment:}
\begin{itemize}
    \item Approximately 1800 machines
    \item Dual-processor 2GHz Intel Xeon with HyperThreading
    \item 4GB memory per machine
    \item Two 160GB IDE disks per machine
    \item Gigabit Ethernet
\end{itemize}

\vspace{0.3cm}

\textbf{Network:}
\begin{itemize}
    \item Two-level tree-shaped switched network
    \item $\approx$ 100-200 Gbps aggregate bandwidth at root
    \item Round-trip time between any pair: $<$ 1 millisecond
\end{itemize}

\vspace{0.3cm}

\textbf{Conditions:}
\begin{itemize}
    \item Weekend afternoon execution
    \item CPUs, disks, network mostly idle
    \item 1-1.5GB reserved for other tasks
\end{itemize}
\end{frame}

\begin{frame}{Benchmark 1: Grep}
\textbf{Task:} Scan through $10^{10}$ 100-byte records, search for rare three-character pattern

\vspace{0.2cm}

\textbf{Configuration:}
\begin{itemize}
    \item Input: 1 TB, split into $\approx$ 64MB pieces ($M = 15,000$)
    \item Output: One file ($R = 1$)
    \item Pattern occurs in 92,337 records
\end{itemize}

\vspace{0.2cm}

\textbf{Results:}
\begin{itemize}
    \item Peak input rate: \textbf{30+ GB/s} (1764 workers)
    \item Computation time: \textbf{150 seconds} (including $\approx$ 60s startup overhead)
    \item Startup overhead due to:
    \begin{itemize}
        \item Program propagation to all workers
        \item GFS interactions to open 1000 input files
        \item Locality optimization information gathering
    \end{itemize}
\end{itemize}

\vspace{0.2cm}

\begin{alertblock}{Key Insight}
Input rate gradually increases as more workers assigned, peaks, then drops as tasks complete
\end{alertblock}
\end{frame}

\begin{frame}{Benchmark 2: Sort (1/2)}
\textbf{Task:} Sort $10^{10}$ 100-byte records ($\approx$ 1 TB of data)

\vspace{0.2cm}

\textbf{Configuration:}
\begin{itemize}
    \item Modeled after TeraSort benchmark
    \item Input split into 64MB pieces ($M = 15,000$)
    \item Output: 4000 files ($R = 4,000$)
    \item Output: 2 TB (2-way replicated GFS files)
\end{itemize}

\vspace{0.2cm}

\textbf{Implementation:}
\begin{itemize}
    \item Map: Extract 10-byte sorting key, emit $\langle$key, line$\rangle$
    \item Reduce: Identity function (passes pairs unchanged)
    \item Less than 50 lines of user code!
\end{itemize}

\vspace{0.2cm}

\textbf{Partitioning:}
\begin{itemize}
    \item Built-in knowledge of key distribution
    \item General sorting: pre-pass to sample keys and compute split-points
\end{itemize}
\end{frame}

\begin{frame}{Benchmark 2: Sort (2/2)}
\textbf{Normal Execution Results:}
\begin{itemize}
    \item \textbf{Total time:} 891 seconds
    \item Input rate peaks: $\approx$ 13 GB/s
    \item Shuffle starts as first map completes
    \item Shuffling done: $\approx$ 600 seconds
    \item Write rate: 2-4 GB/s
    \item All writes finish: $\approx$ 850 seconds
\end{itemize}

\vspace{0.3cm}

\textbf{Observations:}
\begin{itemize}
    \item Input rate $>$ shuffle rate $>$ output rate
    \item Locality optimization: most data read locally
    \item Shuffle rate $>$ output rate: 2 replicas written
    \item Comparable to TeraSort benchmark best result (1057 seconds)
\end{itemize}

\vspace{0.3cm}

\begin{block}{Impact of Backup Tasks}
Without backup tasks: \textbf{1283 seconds} (+44\%) \\
Stragglers significantly impact completion time!
\end{block}
\end{frame}

\begin{frame}{Performance: Machine Failures}
\textbf{Test:} Intentionally killed 200 out of 1746 workers during sort

\vspace{0.3cm}

\textbf{What Happened:}
\begin{itemize}
    \item Underlying cluster scheduler immediately restarted processes
    \item Deaths show up as negative input rate
    \item Previously completed map work disappeared
    \item Re-execution happened relatively quickly
\end{itemize}

\vspace{0.3cm}

\textbf{Results:}
\begin{itemize}
    \item Total time: \textbf{933 seconds}
    \item Increase: Only \textbf{5\%} over normal execution (891s)
\end{itemize}

\vspace{0.5cm}

\begin{exampleblock}{Conclusion}
MapReduce is highly resilient to machine failures! \\
System gracefully handles worker deaths without significant performance degradation.
\end{exampleblock}
\end{frame}

% ============================================
% SECTION 6: EXPERIENCE
% ============================================
\section{Experience at Google}

\begin{frame}{MapReduce at Google}
\textbf{Timeline:}
\begin{itemize}
    \item First version: February 2003
    \item Significant enhancements: August 2003
    \item Locality optimization, dynamic load balancing, etc.
\end{itemize}

\vspace{0.3cm}

\textbf{Growth:}
\begin{itemize}
    \item 0 instances in early 2003
    \item Almost 900 separate instances by September 2004
    \item Used across wide range of domains
\end{itemize}

\vspace{0.3cm}

\textbf{Applications:}
\begin{multicols}{2}
\begin{itemize}
    \item Large-scale machine learning
    \item Clustering problems (Google News, Froogle)
    \item Popular query reports (Google Zeitgeist)
    \item Property extraction from web pages
    \item Large-scale graph computations
    \item Localized search data extraction
\end{itemize}
\end{multicols}
\end{frame}

\begin{frame}{Usage Statistics (August 2004)}
\begin{table}
\centering
\small
\begin{tabular}{lr}
\toprule
\textbf{Metric} & \textbf{Value} \\
\midrule
Number of jobs & 29,423 \\
Average job completion time & 634 secs \\
Machine days used & 79,186 days \\
Input data read & 3,288 TB \\
Intermediate data produced & 758 TB \\
Output data written & 193 TB \\
Average worker machines per job & 157 \\
Average worker deaths per job & 1.2 \\
Average map tasks per job & 3,351 \\
Average reduce tasks per job & 55 \\
Unique map implementations & 395 \\
Unique reduce implementations & 269 \\
Unique map/reduce combinations & 426 \\
\bottomrule
\end{tabular}
\end{table}
\end{frame}

\begin{frame}{Large-Scale Indexing}
\textbf{Most Significant Use:} Complete rewrite of production indexing system

\vspace{0.3cm}

\textbf{Input:} 20+ TB of crawled documents stored in GFS

\vspace{0.3cm}

\textbf{Process:} Sequence of 5-10 MapReduce operations

\vspace{0.3cm}

\textbf{Benefits of Using MapReduce:}

\begin{enumerate}
    \item \textbf{Simpler Code}
    \begin{itemize}
        \item One phase: 3800 lines C++ $\rightarrow$ 700 lines with MapReduce
        \item Fault tolerance, distribution, parallelization hidden in library
    \end{itemize}

    \item \textbf{Better Performance}
    \begin{itemize}
        \item Can keep conceptually unrelated computations separate
        \item No need to mix together to avoid extra data passes
    \end{itemize}

    \item \textbf{Easier to Operate}
    \begin{itemize}
        \item Machine failures, slow machines, networking issues handled automatically
        \item Easy to improve performance: just add more machines!
    \end{itemize}
\end{enumerate}
\end{frame}

\begin{frame}{Why MapReduce is Successful}
\begin{enumerate}
    \item \textbf{Easy to Use}
    \begin{itemize}
        \item No parallel/distributed systems experience needed
        \item Hides complexity of parallelization, fault-tolerance, locality, load balancing
        \item Simple program can run on 1000 machines in half an hour
    \end{itemize}

    \vspace{0.2cm}

    \item \textbf{Widely Applicable}
    \begin{itemize}
        \item Large variety of problems expressible as MapReduce
        \item Used for web search, sorting, data mining, machine learning, and more
    \end{itemize}

    \vspace{0.2cm}

    \item \textbf{Scalable Implementation}
    \begin{itemize}
        \item Scales to thousands of machines
        \item Efficient use of machine resources
        \item Suitable for Google's large computational problems
    \end{itemize}
\end{enumerate}

\vspace{0.3cm}

\begin{alertblock}{Development Impact}
Greatly speeds up development and prototyping cycle \\
Allows programmers to exploit large amounts of resources easily
\end{alertblock}
\end{frame}

% ============================================
% SECTION 7: KEY LEARNINGS
% ============================================
\section{Key Learnings and Conclusions}

\begin{frame}{Key Lessons Learned}
\begin{enumerate}
    \item \textbf{Restricting the Programming Model}
    \begin{itemize}
        \item Makes it easy to parallelize and distribute computations
        \item Enables fault-tolerance through re-execution
        \item Simplicity is powerful!
    \end{itemize}

    \vspace{0.3cm}

    \item \textbf{Network Bandwidth is Scarce}
    \begin{itemize}
        \item Many optimizations target reducing network usage
        \item Locality optimization: read from local disks
        \item Single copy of intermediate data to local disk saves bandwidth
    \end{itemize}

    \vspace{0.3cm}

    \item \textbf{Redundant Execution is Valuable}
    \begin{itemize}
        \item Reduces impact of slow machines (stragglers)
        \item Handles machine failures gracefully
        \item Handles data loss effectively
    \end{itemize}
\end{enumerate}
\end{frame}

\begin{frame}{Impact and Legacy}
\begin{block}{Historical Impact}
This paper (OSDI 2004) became one of the most influential papers in distributed systems and big data processing
\end{block}

\vspace{0.3cm}

\textbf{Inspired Numerous Systems:}
\begin{itemize}
    \item \textbf{Apache Hadoop} (2006) - Open-source MapReduce implementation
    \item \textbf{Apache Spark} - Next-generation data processing
    \item \textbf{Apache Flink, Apache Beam} - Stream processing frameworks
    \item Countless other distributed data processing systems
\end{itemize}

\vspace{0.3cm}

\textbf{Paradigm Shift:}
\begin{itemize}
    \item Democratized large-scale data processing
    \item Made distributed computing accessible to non-experts
    \item Enabled the "Big Data" revolution
    \item Foundation for modern data engineering practices
\end{itemize}
\end{frame}

\begin{frame}{Summary}
\begin{block}{MapReduce in a Nutshell}
A programming model and implementation for processing large datasets on clusters of commodity machines
\end{block}

\vspace{0.3cm}

\textbf{Key Characteristics:}
\begin{itemize}
    \item Simple programming interface (Map and Reduce functions)
    \item Automatic parallelization and distribution
    \item Fault tolerance via re-execution
    \item Locality optimization for performance
    \item Scalable to thousands of machines
\end{itemize}

\vspace{0.3cm}

\textbf{Why It Matters:}
\begin{itemize}
    \item Hides complexity of distributed systems
    \item Allows programmers to focus on logic, not infrastructure
    \item Proven at massive scale at Google
    \item Influenced entire generation of big data systems
\end{itemize}
\end{frame}

\begin{frame}{Questions to Consider}
\begin{enumerate}
    \item What types of problems are \textit{not} well-suited for MapReduce?

    \vspace{0.3cm}

    \item How does MapReduce compare to modern frameworks like Apache Spark?

    \vspace{0.3cm}

    \item What are the trade-offs between simplicity and flexibility in the MapReduce model?

    \vspace{0.3cm}

    \item How would you extend MapReduce to support iterative algorithms?

    \vspace{0.3cm}

    \item What role does MapReduce play in today's data processing landscape?
\end{enumerate}
\end{frame}

\begin{frame}[standout]
\Huge Thank You!

\vspace{1cm}

\normalsize
Questions?

\vspace{1cm}

\footnotesize
Paper: Dean, Jeffrey, and Sanjay Ghemawat. \\
"MapReduce: Simplified data processing on large clusters." \\
\textit{OSDI 2004}
\end{frame}

\appendix

\begin{frame}[fragile]{Appendix: Word Count Full Implementation}
\begin{lstlisting}[language=C++, basicstyle=\ttfamily\tiny]
#include "mapreduce/mapreduce.h"

// User's map function
class WordCounter : public Mapper {
  public:
    virtual void Map(const MapInput& input) {
      const string& text = input.value();
      const int n = text.size();
      for (int i = 0; i < n; ) {
        while ((i < n) && isspace(text[i])) i++;
        int start = i;
        while ((i < n) && !isspace(text[i])) i++;
        if (start < i)
          Emit(text.substr(start,i-start),"1");
      }
    }
};
REGISTER_MAPPER(WordCounter);

// User's reduce function
class Adder : public Reducer {
  virtual void Reduce(ReduceInput* input) {
    int64 value = 0;
    while (!input->done()) {
      value += StringToInt(input->value());
      input->NextValue();
    }
    Emit(IntToString(value));
  }
};
REGISTER_REDUCER(Adder);
\end{lstlisting}
\end{frame}

\begin{frame}{Appendix: Related Work}
\textbf{Parallel Programming Models:}
\begin{itemize}
    \item Bulk Synchronous Programming (BSP)
    \item MPI primitives
    \item Parallel prefix computations
\end{itemize}

\vspace{0.2cm}

\textbf{Distributed Systems:}
\begin{itemize}
    \item Active Disks - computation near storage
    \item Charlotte System - eager scheduling
    \item Condor - cluster management
\end{itemize}

\vspace{0.2cm}

\textbf{Sorting and Data Processing:}
\begin{itemize}
    \item NOW-Sort - distributed sorting
    \item River - distributed queues
    \item BAD-FS - wide-area network execution
\end{itemize}

\vspace{0.2cm}

\textbf{MapReduce Differences:}
\begin{itemize}
    \item Simpler programming model
    \item Automatic fault tolerance at scale
    \item Focus on commodity hardware
    \item Proven at production scale
\end{itemize}
\end{frame}

\end{document}
